% Standalone problem document, generated from the adjacent README.md.
% Compile with XeLaTeX. Mathematical content is maintained in Markdown.
\documentclass[11pt,a4paper]{article}
\usepackage[fontsize=10.5pt]{scrextend}
\usepackage[margin=22mm,headheight=15pt,headsep=9mm,footskip=11mm]{geometry}
\usepackage{fontspec}
\setmainfont{texgyrepagella}[Extension=.otf,UprightFont=*-regular,BoldFont=*-bold,ItalicFont=*-italic,BoldItalicFont=*-bolditalic]
\setsansfont{texgyreheros}[Extension=.otf,UprightFont=*-regular,BoldFont=*-bold,ItalicFont=*-italic,BoldItalicFont=*-bolditalic]
\setmonofont{texgyrecursor}[Extension=.otf,UprightFont=*-regular,BoldFont=*-bold,ItalicFont=*-italic,BoldItalicFont=*-bolditalic,Scale=MatchLowercase]
\usepackage{amsmath,amssymb}
\usepackage{unicode-math}
\setmathfont{texgyrepagella-math.otf}
\usepackage{xcolor,microtype,enumitem,needspace,fancyhdr}
\usepackage{bookmark}
\definecolor{ink}{HTML}{17344B}
\definecolor{accent}{HTML}{197D80}
\definecolor{muted}{HTML}{596773}
\hypersetup{colorlinks=true,linkcolor=accent,urlcolor=accent,pdftitle={TR-03 - Sharp
gap between volume sampling and the worst matrix with a prescribed
spectrum},pdfauthor={Open Problems in Numerical Linear Algebra}}
\urlstyle{same}
\setlength{\parindent}{0pt}
\setlength{\parskip}{5pt plus 1pt minus 1pt}
\linespread{1.04}
\setlength{\emergencystretch}{3em}
\widowpenalty=10000
\clubpenalty=10000
\displaywidowpenalty=10000
\allowdisplaybreaks[1]
\setlist{nosep,leftmargin=1.5em}
\providecommand{\tightlist}{\setlength{\itemsep}{0pt}\setlength{\parskip}{0pt}}
\providecommand{\pandocbounded}[1]{#1}
\pagestyle{fancy}
\fancyhf{}
\fancyhead[L]{\sffamily\footnotesize\color{muted}OPEN PROBLEMS IN NUMERICAL LINEAR ALGEBRA}
\fancyhead[R]{\sffamily\bfseries\footnotesize\color{accent}TR-03}
\fancyfoot[L]{\parbox[b]{0.9\textwidth}{\sffamily\fontsize{8}{10}\selectfont\color{muted}Editorial ratings; a bounded search cannot certify that a problem remains unsolved.\\Literature check: 2026-09-10}}
\fancyfoot[R]{\sffamily\footnotesize\color{muted}\thepage}
\renewcommand{\headrulewidth}{0.3pt}
\renewcommand{\headrule}{\hbox to\headwidth{\color{accent}\leaders\hrule height \headrulewidth\hfill}}
\makeatletter
\renewcommand{\section}{\Needspace{5\baselineskip}\@startsection{section}{1}{0pt}{12pt plus 2pt minus 2pt}{4pt}{\sffamily\bfseries\large\color{ink}}}
\renewcommand{\subsection}{\Needspace{4\baselineskip}\@startsection{subsection}{2}{0pt}{9pt plus 1pt minus 1pt}{3pt}{\sffamily\bfseries\normalsize\color{ink}}}
\makeatother
\setcounter{secnumdepth}{0}
\begin{document}
{\sffamily\small\color{muted}Randomized and low-rank approximation\par}
\vspace{3pt}
{\raggedright\hyphenpenalty=10000\exhyphenpenalty=10000\sffamily\bfseries\fontsize{21}{24}\selectfont\color{ink}Sharp
gap between volume sampling and the worst matrix with a prescribed
spectrum\par}
\vspace{7pt}
{\sffamily\small\textbf{Difficulty:} challenging\quad\textbf{Importance:} interesting
to the community\par}
{\sffamily\small\bfseries\color{accent}Status: Open\par}
\vspace{4pt}
{\color{accent}\hrule height 0.8pt}
\vspace{6pt}
\textbf{Rating rationale:} Challenging because the adversarial
eigenvector choice and optimal subset must be compared sharply;
community impact is spectrum-sensitive Nyström and column selection.

\subsection{Problem statement}\label{problem-statement}

For \(n\ge3\), \(1\le k\le n-2\), and \(\lambda\in(0,\infty)^n\), write
\(K_V=V^T\operatorname{diag}(\lambda)V\), \(V\in O(n)\), and define

\[
x_k(\lambda)=\max_{V\in O(n)}\min_{|I|=k}
\operatorname{tr}\bigl(K_V-(K_V)_{:I}(K_V)_{II}^{-1}(K_V)_{I:}\bigr),
\quad
y_k(\lambda)=(k+1)\frac{e_{k+1}(\lambda)}{e_k(\lambda)},
\]

where \(e_j(\lambda)=\sum_{|I|=j}\prod_{i\in I}\lambda_i\). Determine,
up to universal multiplicative constants, the dependence on \(n,k\) of

\[
R_{n,k}=\sup_{\lambda\in(0,\infty)^n}\frac{y_k(\lambda)}{x_k(\lambda)}.
\]

Here \(y_k\) is the expected trace error when the selected subset has
probability \(\det((K_V)_{II})/e_k(\lambda)\). The question compares
that expectation with optimal subset selection after an adversary
chooses the eigenvectors. This order of quantifiers is essential. The
endpoint \(R_{n,n-1}=1\) is known and excluded. The displayed supremum
is a concrete subquestion of the source's request for spectrum-dependent
tightness bounds; stronger bounds retaining the full spectrum would also
be valuable.

\subsection{References}\label{references}

Amsel et al., \href{https://arxiv.org/html/2602.05394v3}{workshop
report}, §4.2, Problem 4.7 and equations (15)--(16). Mark Fornace and
Michael Lindsey, \href{https://arxiv.org/html/2407.01698v2}{\emph{Column
and Row Subset Selection Using Nuclear Scores: Algorithms and Theory for
Nyström Approximation, CUR Decomposition, and Graph Laplacian
Reduction}}, §4 and Appendix C, for determinantal sampling and the
elementary-symmetric-polynomial formulas.

\subsection{Status check}\label{status-check}

Searches included \textit{"minimax" "volume sampling" "2026"},
\textit{"volume sampling" "worst" "spectrum" "2026"}, and
\textit{"volume sampling" "tightness" Fornace}. No matching sharp
estimate for \(R_{n,k}\) was located. Generic column-subset
approximation guarantees do not by themselves settle this spectral
minimax ratio.

\subsection{Audit --- 2026-09-10}\label{audit-2026-09-10}

Rechecked \href{https://arxiv.org/html/2602.05394v3}{workshop Problem
4.7 and equations (15)--(16)}. They support the stated spectral-minimax
subquestion. Volume-sampling, worst-spectrum, and tightness searches
found no sharp joint estimate for the displayed ratio; generic subset
guarantees leave its quantifier order unresolved.
\end{document}
